\section{Discussion}
\label{sec:discussion}

Here we interpret the classification and XAI results from Section~\ref{sec:results}, compare them to earlier studies. And plainly lay out what this work proves and what it does not.

\subsection{Interpretation of results}

The primary model hits 93.05\% accuracy and 0.9965 macro-AUC (Table~\ref{tbl:3}), and a three-seed replication confirms it ($0.9210\pm0.0131$). Unlike the single-centre GBCU dataset studied before, the 6.6-point F1 difference---between Cholecystitis (0.9601) and Perforation (0.8944)---doesn't follow training-set size. It follows instead the visual overlap between Perforation and Gallstones, and between Perforation and Membranous/Gangrenous Cholecystitis (Section~\ref{sec:results}.2).

Fig.~\ref{fig:4a}'s confusion matrix reveals the biggest mistake: Polyps~\&~Cholesterol~Crystals~$\leftrightarrow$~Adenomyomatosis, a benign-to-benign swap. A more urgent slip is Carcinoma~$\rightarrow$~Gallstones (12 cases), a missed-malignancy pattern, but it's still a minority.

When we look at cross-replicate consistency, Grad-CAM++ and Score-CAM come out on top. They beat Grad-CAM, Saliency Maps, and Eigen-CAM (Section~\ref{sec:results}.3). The result: sharper, more class-discriminative heatmaps, and the highest, most stable consistency ($\bar{O}^c=0.944$ and $0.937$ versus 0.769, 0.242, and 0.681 for Grad-CAM, Eigen-CAM, and Saliency).

But faithfulness tells a different story. Every CAM family method, even plain Grad-CAM, outperforms the random control (Section~\ref{sec:results}.4). On the deletion/insertion metric, plain Grad-CAM shows the biggest gap: 0.480 against 0.473 and 0.456. The difference is so narrow that the authors avoid ranking the methods that way.

Sharpness and reproducibility lean toward Grad-CAM++ and Score-CAM, yet faithfulness is level across the trio. No method wins on every metric, so picking one depends on what you value.

Eigen-CAM's shaky consistency ($\bar{O}^c=0.242\pm0.343$, dropping below 0.10 for three classes) is predictable. It's class-agnostic by design, so its target differs from the four class-conditional methods. We keep it as a cheap, gradient-free baseline, not a deployment tool. Its lower consistency alone shouldn't signal poorer explanations.

The six per-class observations below are the authors' own unblinded visual read of the heatmaps, offered as hypotheses about what the model may be attending to, not as clinical findings. Neither author has clinical/radiological training, and none of these observations have been checked against an independent clinician's judgement or region-of-interest annotation; this paper's own position throughout is that visual plausibility is an insufficient reliability standard on its own, so these should be read as candidates for the radiologist validation protocol below (Section~\ref{sec:discussion}.3), not as evidence in themselves.

\textbf{Gallstones.} Grad-CAM++ heatmaps consistently mark the region below the gallbladder that, on visual inspection, resembles the posterior acoustic shadow conventionally associated with cholelithiasis. Whether the model has learned a clinically valid feature, as opposed to a strong but over-generalised correlate, is not something this observation can distinguish (Section~\ref{sec:limitations}). Consistent with over-generalisation, Gallstones also receives the most confusion-pair traffic: both Perforation and Carcinoma most often fall into it (Section~\ref{sec:results}.2).

\textbf{Cholecystitis and Membranous/Gangrenous Cholecystitis.} Both heatmaps visually concentrate on the gallbladder wall; the Membranous/Gangrenous maps appear more diffuse across the wall, a pattern we note without asserting it reflects the more extensive mural involvement of that stage. The two classes are mutually confused in 12 cases (Section~\ref{sec:results}.2), consistent with them lying on a severity spectrum rather than being visually independent categories.

\textbf{Perforation.} Heatmaps for this class visually concentrate on wall-defect and pericholecystic regions, though it remains the weakest class by F1 (0.8944, Table~\ref{tbl:4}). Its biggest mistakes are to Gallstones (17 cases) and to Membranous/Gangrenous Cholecystitis (11); the failure-case heatmap in Section~\ref{sec:results}.2 shows a true perforation read as Gallstones (99.6\%) from what visually appears to be a shadow-adjacent structure rather than a clear defect.

\textbf{Carcinoma.} Grad-CAM++ heatmaps visually concentrate on a focal wall area; we note this without asserting a correspondence to the irregular thickening and mass-like lesions that characterise GBC clinically, since that correspondence has not been independently checked. The Gallstones-directed mistakes (12 cases, Section~\ref{sec:results}.2) remain the most consequential confusion pattern in this dataset on face value: a missed malignancy read as benign.

\textbf{Polyps \& Cholesterol Crystals and Adenomyomatosis.} These form the largest confusion pair (25 cases, Section~\ref{sec:results}.2), and their heatmaps visually overlap: both concentrate on small, focal mural spots rather than a diffuse wall region.

\textbf{Abdomen \& Retroperitoneum and Wall Thickening.} These broader, less gallbladder-specific classes show more dispersed heatmap attention. Grad-CAM++ appears visually sharper than plain Grad-CAM here (Table~\ref{tbl:6}: Grad-CAM falls to 0.595 for Abdomen \& Retroperitoneum and 0.630 for Wall Thickening), though whether that sharper attention picks out anything clinically meaningful, rather than simply differing in texture, is exactly the kind of question the planned radiologist review (Section~\ref{sec:discussion}.3) and not further author inspection needs to answer.

\subsection{Reliability and trustworthiness of explanations}

Sections~\ref{sec:results}.3 and \ref{sec:results}.4 build a cumulative case for trustworthiness along five axes. The explanations replicate across independently trained copies. This is most evident with Grad-CAM++ and Score-CAM. They mirror what the model actually uses. They stay stable under transforms that a diagnosis should ignore---brightness is the cleanest result. And they depend genuinely on learned weights. On the margin-based attribution evidence specifically, there is no clear shortcut driving them. The crop-ablation leg of that test isn't fully conclusive (Section~\ref{sec:results}.4). Passing these five axes shows the explanations faithfully, reproducibly, and stably describe what the model does. That alone, however, does not prove clinical trustworthiness. A classifier can be internally consistent, well-calibrated, and give faithful explanations for a process that has never been validated against clinician ground truth. The margin-based shortcut audit only checks for shortcuts tied to the image margin. It does not rule out other, unmeasured shortcuts a classifier could learn. In sum, this paper's XAI-method comparison is measured and reproducible. But its classifier's explanations of disease have not yet been validated against an independent clinician's judgement (Section~\ref{sec:future-research}).

\subsection{Comparison with previous studies and practical implications}

Section~\ref{sec:related-work} reviewed the gallbladder-ultrasound and medical-XAI literature in detail. This section summarises how this paper's results position it against that work. Table~\ref{tbl:9} makes the comparison concrete. Prior gallbladder-ultrasound classifiers \cite{ref7}, \cite{ref17}, \cite{ref18}, \cite{ref19}, \cite{ref22} have delivered a classifier, a dataset, or, in RadFormer's case \cite{ref19}, quantitative localization for a single explanation mechanism. Nabil et~al. \cite{ref57} and Kumar et~al. \cite{ref58} each used two XAI methods together, but compared them qualitatively. No source identified in this search combines five XAI methods, a quantitative cross-training consistency metric, and grounded failure-case analysis for gallbladder ultrasound. That is the specific, narrower claim this paper makes. It is not the first to apply XAI or any quantitative XAI evaluation to this domain. Nor is it the broader claim of accuracy gains reported by concurrent anatomy-aware architectures on GBCU \cite{ref22}, \cite{ref23}. No claim is made here to advance classification accuracy. The goal is to establish, with more measurement than is typical in this literature, whether the explanations of a standard architecture can be trusted.

\begin{table*}[!ht]
\centering
\caption{Comparison with Prior Gallbladder-Ultrasound Classification and XAI Work Only. General-purpose method papers (Grad-CAM, Grad-CAM++, Score-CAM's original sources) are excluded, since they never claimed to address this domain and their inclusion would inflate the apparent novelty gap rather than compare like with like. ``XAI Methods'' counts distinct post-hoc methods used, regardless of whether their comparison was qualitative or quantitative; the next column states which explicitly.}
\label{tbl:9}
\begin{tabular}{lccc}
\toprule
Reference & XAI Methods (qual./quant.) & Quant. Consistency Metric & Grounded Failure-Case Analysis \\
\midrule
Turki et al.~\cite{ref17} (dataset release) & 0 & $\times$ & $\times$ \\
Obaid et al.~\cite{ref18} & 1 (qual.) & $\times$ & $\times$ \\
Nabil et al.~\cite{ref57} & 2 (qual.) & $\times$ & $\times$ \\
Kumar et al.~\cite{ref58} & 2 (qual.\ hybrid) & $\times$ & $\times$ \\
Basu et al.~\cite{ref7} (GBCU, separate dataset) & 0 & $\times$ & $\times$ \\
RadFormer~\cite{ref19} (GBCU, separate dataset) & 1 (quant.\ localization) & $\times$ & $\times$ \\
Hasan et al.~\cite{ref22} (GBCU, separate dataset) & 1 (qual.) & $\times$ & $\times$ \\
\textbf{Ours (UIdataGB)} & \textbf{5 (quant.)} & \textbf{\checkmark} & \textbf{\checkmark} \\
\bottomrule
\end{tabular}

\vspace{2pt}
\raggedright\footnotesize All rows are gallbladder-focused studies, either on UIdataGB directly (this paper's dataset) or on the separate GBCU dataset (noted per row); rows report what each study actually did, not what it did not do. Turki et al. is the UIdataGB dataset-release paper used in this study; Obaid et al. is the closest prior classifier, trained on the same raw 10,692-image release. Nabil et al.\ and Kumar et al.\ are the closest prior XAI comparators, each combining Grad-CAM with LIME on gallbladder-ultrasound data, but reporting the comparison qualitatively rather than against the reliability metrics (consistency, faithfulness, stability, sanity check) used here. RadFormer provides the strongest prior quantitative XAI evidence in this domain, via attention-map localization against radiologist annotations on GBCU, for a single explanation mechanism rather than a multi-method comparison. No row other than ``Ours'' combines $\ge$5 XAI methods with a quantitative cross-training consistency metric and grounded failure-case analysis for gallbladder ultrasound; this is a narrower claim than being the first application of XAI, or of any quantitative XAI evaluation, to this domain, and it should be read alongside this paper's own Section~\ref{sec:results}.5 correction, since the quantitative consistency claim in the ``Ours'' row above was established on the Stage~1 split and only partially re-verified after the leakage correction.
\end{table*}

\textbf{What this paper does and does not support.} The consistency, faithfulness, stability, sanity-check, calibration. And shortcut results (Sections~\ref{sec:results}.3 and \ref{sec:results}.4) are evidence that the model's explanations are reproducible, faithful to its decision, stable, weight-dependent, and free of obvious margin shortcuts. They also show that its probabilities can be calibrated. They are not yet evidence that the model is ready for decision support. That would require quantitative localization of the explanation against clinician-defined regions of interest, plus formal clinician review of the heatmaps. Neither can be produced without expert annotation (Section~\ref{sec:future-research}).

\textbf{Monitoring explanation reproducibility.} The consistency score $O_r^c$ could serve as an ongoing quality check when a model is retrained on updated data from this same release. If $O_r^c$ drops below $\theta=0.70$ for a class after retraining, this flags that the update materially changed what the model attends to for that class. A targeted review would then be triggered before the updated model replaces the deployed one.

\textbf{Radiologist validation protocol (planned).} Before these findings can support any clinical claim, a structured clinician evaluation is required. It follows the protocol in the roadmap (Section~\ref{sec:future-research}): (1)~two or more radiologists independently rate heatmap clinical relevance and whether the highlighted region supports the predicted class, on a 5-point Likert scale, with a separate binary flag for attention on non-clinical artefacts (text, calipers, borders). (2)~Inter-rater agreement is reported via Cohen's $\kappa$. (3)~Model attention is compared against radiologist-drawn gallbladder and pathology bounding boxes using IoU, Dice, and pointing-game accuracy. The stratified per-class Grad-CAM++ gallery already generated for this study (Supplementary Material A, 54 images across all nine classes) is a natural starting panel for such a rating protocol once clinician participation is secured.
